\section{万历朝的补充编年节点（二）}
以下节点继续沿同一编年脉络展开。每一锚点均紧接可供核查的事件叙述；来源能确认何种行动、文书或变动，不能自动消除不同记录之间的立场差异。

\subsection{1606年前后的续接}

\eventanchor{ML-1606-002}{明·1606（万历三十四年）}
\paragraph{1606（万历三十四年）：明式步枪均附有插入式刺刀。} 这项记录把本章所见的制度运作落实到具体年份。其形成背景是：继承、官僚协调或皇权运作的压力使问题进入朝廷程序。 它带来的、而且仍可由后续材料追踪的改变是：它影响后续人事与政策议程；动机和派系标签不能由单一叙事裁定。 《神宗实录》万历三十四年条；《明史·本纪》及相关列传；诏令、奏疏或会典版本 提供相互检验的起点；但译写候选以编年材料定位；实录记载反映朝廷选择，崇祯期须另以奏疏、地方及敌对政权材料交叉。 因此本段仅据材料确认事件的时序、行动和可见后果，并把政策文本与地方实践、报称数字和社会经历分开处理。

\eventanchor{ML-1606-003}{明·1606（万历三十四年）}
\paragraph{1606（万历三十四年）：后金兴起、朝鲜交涉和海上网络的本年多方材料揭示区域关系变化，政权名称与译名须按原文说明。（材料核查重点：诏令或奏议的形成与发受链）} 这项记录把本章所见的制度运作落实到具体年份。其形成背景是：朝贡名义、边境安全与港口贸易的利益使交涉持续发生。 它带来的、而且仍可由后续材料追踪的改变是：它为后续册封、互市或海防安排留下文书与跨政权比较线索。 《神宗实录》万历三十四年条；《明史·外国传》；《朝鲜王朝实录》、琉球《历代宝案》或海外档案 提供相互检验的起点；但桥接条目只标示可回查的年度问题与材料链；实录有中央选择性，崇祯期尤其需要奏疏、地方、朝鲜及满文材料交叉。；本条特别核对诏令或奏议的形成与发受链。 因此本段仅据材料确认事件的时序、行动和可见后果，并把政策文本与地方实践、报称数字和社会经历分开处理。

\eventanchor{ML-1606-004}{明·1606（万历三十四年）}
\paragraph{1606（万历三十四年）：本年灾害、疫病或地方赈济线索应以受灾地区文书和地方志复核，中央奏报不能覆盖全部社会。} 这项记录把本章所见的制度运作落实到具体年份。其形成背景是：自然风险与地方报告促成朝廷或地方的应对记录。 它带来的、而且仍可由后续材料追踪的改变是：赈济、迁徙和损失的实际范围仍须以受灾地区材料分别核验。 《神宗实录》万历三十四年条；《明史·五行志》；受灾府县地方志与奏疏 提供相互检验的起点；但桥接条目只标示可回查的年度问题与材料链；实录有中央选择性，崇祯期尤其需要奏疏、地方、朝鲜及满文材料交叉。 因此本段仅据材料确认事件的时序、行动和可见后果，并把政策文本与地方实践、报称数字和社会经历分开处理。

\eventanchor{ML-1606-005}{明·1606（万历三十四年）}
\paragraph{1606（万历三十四年）：后金兴起、朝鲜交涉和海上网络的本年多方材料揭示区域关系变化，政权名称与译名须按原文说明。（材料核查重点：报称信息的来源、抄传与行政分类）} 这项记录把本章所见的制度运作落实到具体年份。其形成背景是：朝贡名义、边境安全与港口贸易的利益使交涉持续发生。 它带来的、而且仍可由后续材料追踪的改变是：它为后续册封、互市或海防安排留下文书与跨政权比较线索。 《神宗实录》万历三十四年条；《明史·外国传》；《朝鲜王朝实录》、琉球《历代宝案》或海外档案 提供相互检验的起点；但桥接条目只标示可回查的年度问题与材料链；实录有中央选择性，崇祯期尤其需要奏疏、地方、朝鲜及满文材料交叉。；本条特别核对报称信息的来源、抄传与行政分类。 因此本段仅据材料确认事件的时序、行动和可见后果，并把政策文本与地方实践、报称数字和社会经历分开处理。

\eventanchor{ML-1606-006}{明·1606（万历三十四年）}
\paragraph{1606（万历三十四年）：辽东战事、边镇防务和西南兵事的本年军令记录资源调动，战报数字应与朝鲜、满文或地方材料比照。} 这项记录把本章所见的制度运作落实到具体年份。其形成背景是：前期边防、地方武装或跨区域资源调动的压力推动了该行动。 它带来的、而且仍可由后续材料追踪的改变是：它改变了后续调兵、补给或地方秩序的条件；战果和伤亡须分开核对。 《神宗实录》万历三十四年条；《明史·兵志》及相关列传；边镇、地方志或对方材料 提供相互检验的起点；但桥接条目只标示可回查的年度问题与材料链；实录有中央选择性，崇祯期尤其需要奏疏、地方、朝鲜及满文材料交叉。 因此本段仅据材料确认事件的时序、行动和可见后果，并把政策文本与地方实践、报称数字和社会经历分开处理。

\eventanchor{ML-1607-001}{明·1607（万历三十五年）}
\paragraph{1607（万历三十五年）：欧几里得几何原理前六册被翻译成中文} 这项记录把本章所见的制度运作落实到具体年份。其形成背景是：制度、教育或知识传播的需要推动了相关行动或文本形成。 它带来的、而且仍可由后续材料追踪的改变是：它提供制度和传播史的锚点，不能据此推断社会整体接受。 《神宗实录》万历三十五年条；《明史·选举志》或《艺文志》；文集、碑刻或书目材料 提供相互检验的起点；但译写候选以编年材料定位；实录记载反映朝廷选择，崇祯期须另以奏疏、地方及敌对政权材料交叉。 因此本段仅据材料确认事件的时序、行动和可见后果，并把政策文本与地方实践、报称数字和社会经历分开处理。

\eventanchor{ML-1607-002}{明·1607（万历三十五年）}
\paragraph{1607（万历三十五年）：本年灾害、疫病或地方赈济线索应以受灾地区文书和地方志复核，中央奏报不能覆盖全部社会。} 这项记录把本章所见的制度运作落实到具体年份。其形成背景是：自然风险与地方报告促成朝廷或地方的应对记录。 它带来的、而且仍可由后续材料追踪的改变是：赈济、迁徙和损失的实际范围仍须以受灾地区材料分别核验。 《神宗实录》万历三十五年条；《明史·五行志》；受灾府县地方志与奏疏 提供相互检验的起点；但桥接条目只标示可回查的年度问题与材料链；实录有中央选择性，崇祯期尤其需要奏疏、地方、朝鲜及满文材料交叉。 因此本段仅据材料确认事件的时序、行动和可见后果，并把政策文本与地方实践、报称数字和社会经历分开处理。

\eventanchor{ML-1607-003}{明·1607（万历三十五年）}
\paragraph{1607（万历三十五年）：东林、书院、刻书和宗教活动的本年材料记录精英网络，社团存在不等于一致政治立场。（材料核查重点：诏令或奏议的形成与发受链）} 这项记录把本章所见的制度运作落实到具体年份。其形成背景是：制度、教育或知识传播的需要推动了相关行动或文本形成。 它带来的、而且仍可由后续材料追踪的改变是：它提供制度和传播史的锚点，不能据此推断社会整体接受。 《神宗实录》万历三十五年条；《明史·选举志》或《艺文志》；文集、碑刻或书目材料 提供相互检验的起点；但桥接条目只标示可回查的年度问题与材料链；实录有中央选择性，崇祯期尤其需要奏疏、地方、朝鲜及满文材料交叉。；本条特别核对诏令或奏议的形成与发受链。 因此本段仅据材料确认事件的时序、行动和可见后果，并把政策文本与地方实践、报称数字和社会经历分开处理。

\eventanchor{ML-1607-004}{明·1607（万历三十五年）}
\paragraph{1607（万历三十五年）：万历末至天启初的立储、内阁和言路材料标记中枢运作的堵塞与调整，派系归类需回到具体文书。} 这项记录把本章所见的制度运作落实到具体年份。其形成背景是：继承、官僚协调或皇权运作的压力使问题进入朝廷程序。 它带来的、而且仍可由后续材料追踪的改变是：它影响后续人事与政策议程；动机和派系标签不能由单一叙事裁定。 《神宗实录》万历三十五年条；《明史·本纪》及相关列传；诏令、奏疏或会典版本 提供相互检验的起点；但桥接条目只标示可回查的年度问题与材料链；实录有中央选择性，崇祯期尤其需要奏疏、地方、朝鲜及满文材料交叉。 因此本段仅据材料确认事件的时序、行动和可见后果，并把政策文本与地方实践、报称数字和社会经历分开处理。

\eventanchor{ML-1607-005}{明·1607（万历三十五年）}
\paragraph{1607（万历三十五年）：东林、书院、刻书和宗教活动的本年材料记录精英网络，社团存在不等于一致政治立场。（材料核查重点：报称信息的来源、抄传与行政分类）} 这项记录把本章所见的制度运作落实到具体年份。其形成背景是：制度、教育或知识传播的需要推动了相关行动或文本形成。 它带来的、而且仍可由后续材料追踪的改变是：它提供制度和传播史的锚点，不能据此推断社会整体接受。 《神宗实录》万历三十五年条；《明史·选举志》或《艺文志》；文集、碑刻或书目材料 提供相互检验的起点；但桥接条目只标示可回查的年度问题与材料链；实录有中央选择性，崇祯期尤其需要奏疏、地方、朝鲜及满文材料交叉。；本条特别核对报称信息的来源、抄传与行政分类。 因此本段仅据材料确认事件的时序、行动和可见后果，并把政策文本与地方实践、报称数字和社会经历分开处理。

\eventanchor{ML-1607-006}{明·1607（万历三十五年）}
\paragraph{1607（万历三十五年）：后金兴起、朝鲜交涉和海上网络的本年多方材料揭示区域关系变化，政权名称与译名须按原文说明。} 这项记录把本章所见的制度运作落实到具体年份。其形成背景是：朝贡名义、边境安全与港口贸易的利益使交涉持续发生。 它带来的、而且仍可由后续材料追踪的改变是：它为后续册封、互市或海防安排留下文书与跨政权比较线索。 《神宗实录》万历三十五年条；《明史·外国传》；《朝鲜王朝实录》、琉球《历代宝案》或海外档案 提供相互检验的起点；但桥接条目只标示可回查的年度问题与材料链；实录有中央选择性，崇祯期尤其需要奏疏、地方、朝鲜及满文材料交叉。 因此本段仅据材料确认事件的时序、行动和可见后果，并把政策文本与地方实践、报称数字和社会经历分开处理。

\eventanchor{ML-1608-001}{明·1608（万历三十六年）}
\paragraph{1608（万历三十六年）：矿税、加派、漕运和辽饷的本年奏疏与账目提供财政压力线索，预算、欠额和实解不得混同。（材料核查重点：诏令或奏议的形成与发受链）} 这项记录把本章所见的制度运作落实到具体年份。其形成背景是：财政征解、交通工程或货币支付的压力使相关措施进入政务。 它带来的、而且仍可由后续材料追踪的改变是：其法定安排不等于地方执行，后续须以地域材料追踪负担与成效。 《神宗实录》万历三十六年条；《明会典》相关条目；地方志、黄册/契约/河工材料 提供相互检验的起点；但桥接条目只标示可回查的年度问题与材料链；实录有中央选择性，崇祯期尤其需要奏疏、地方、朝鲜及满文材料交叉。；本条特别核对诏令或奏议的形成与发受链。 因此本段仅据材料确认事件的时序、行动和可见后果，并把政策文本与地方实践、报称数字和社会经历分开处理。

\eventanchor{ML-1608-002}{明·1608（万历三十六年）}
\paragraph{1608（万历三十六年）：万历末至天启初的立储、内阁和言路材料标记中枢运作的堵塞与调整，派系归类需回到具体文书。} 这项记录把本章所见的制度运作落实到具体年份。其形成背景是：继承、官僚协调或皇权运作的压力使问题进入朝廷程序。 它带来的、而且仍可由后续材料追踪的改变是：它影响后续人事与政策议程；动机和派系标签不能由单一叙事裁定。 《神宗实录》万历三十六年条；《明史·本纪》及相关列传；诏令、奏疏或会典版本 提供相互检验的起点；但桥接条目只标示可回查的年度问题与材料链；实录有中央选择性，崇祯期尤其需要奏疏、地方、朝鲜及满文材料交叉。 因此本段仅据材料确认事件的时序、行动和可见后果，并把政策文本与地方实践、报称数字和社会经历分开处理。

\eventanchor{ML-1608-003}{明·1608（万历三十六年）}
\paragraph{1608（万历三十六年）：本年灾害、疫病或地方赈济线索应以受灾地区文书和地方志复核，中央奏报不能覆盖全部社会。（材料核查重点：诏令或奏议的形成与发受链）} 这项记录把本章所见的制度运作落实到具体年份。其形成背景是：自然风险与地方报告促成朝廷或地方的应对记录。 它带来的、而且仍可由后续材料追踪的改变是：赈济、迁徙和损失的实际范围仍须以受灾地区材料分别核验。 《神宗实录》万历三十六年条；《明史·五行志》；受灾府县地方志与奏疏 提供相互检验的起点；但桥接条目只标示可回查的年度问题与材料链；实录有中央选择性，崇祯期尤其需要奏疏、地方、朝鲜及满文材料交叉。；本条特别核对诏令或奏议的形成与发受链。 因此本段仅据材料确认事件的时序、行动和可见后果，并把政策文本与地方实践、报称数字和社会经历分开处理。

\eventanchor{ML-1608-004}{明·1608（万历三十六年）}
\paragraph{1608（万历三十六年）：矿税、加派、漕运和辽饷的本年奏疏与账目提供财政压力线索，预算、欠额和实解不得混同。（材料核查重点：报称信息的来源、抄传与行政分类）} 这项记录把本章所见的制度运作落实到具体年份。其形成背景是：财政征解、交通工程或货币支付的压力使相关措施进入政务。 它带来的、而且仍可由后续材料追踪的改变是：其法定安排不等于地方执行，后续须以地域材料追踪负担与成效。 《神宗实录》万历三十六年条；《明会典》相关条目；地方志、黄册/契约/河工材料 提供相互检验的起点；但桥接条目只标示可回查的年度问题与材料链；实录有中央选择性，崇祯期尤其需要奏疏、地方、朝鲜及满文材料交叉。；本条特别核对报称信息的来源、抄传与行政分类。 因此本段仅据材料确认事件的时序、行动和可见后果，并把政策文本与地方实践、报称数字和社会经历分开处理。

\eventanchor{ML-1608-005}{明·1608（万历三十六年）}
\paragraph{1608（万历三十六年）：本年灾害、疫病或地方赈济线索应以受灾地区文书和地方志复核，中央奏报不能覆盖全部社会。（材料核查重点：报称信息的来源、抄传与行政分类）} 这项记录把本章所见的制度运作落实到具体年份。其形成背景是：自然风险与地方报告促成朝廷或地方的应对记录。 它带来的、而且仍可由后续材料追踪的改变是：赈济、迁徙和损失的实际范围仍须以受灾地区材料分别核验。 《神宗实录》万历三十六年条；《明史·五行志》；受灾府县地方志与奏疏 提供相互检验的起点；但桥接条目只标示可回查的年度问题与材料链；实录有中央选择性，崇祯期尤其需要奏疏、地方、朝鲜及满文材料交叉。；本条特别核对报称信息的来源、抄传与行政分类。 因此本段仅据材料确认事件的时序、行动和可见后果，并把政策文本与地方实践、报称数字和社会经历分开处理。

\eventanchor{ML-1608-006}{明·1608（万历三十六年）}
\paragraph{1608（万历三十六年）：东林、书院、刻书和宗教活动的本年材料记录精英网络，社团存在不等于一致政治立场。} 这项记录把本章所见的制度运作落实到具体年份。其形成背景是：制度、教育或知识传播的需要推动了相关行动或文本形成。 它带来的、而且仍可由后续材料追踪的改变是：它提供制度和传播史的锚点，不能据此推断社会整体接受。 《神宗实录》万历三十六年条；《明史·选举志》或《艺文志》；文集、碑刻或书目材料 提供相互检验的起点；但桥接条目只标示可回查的年度问题与材料链；实录有中央选择性，崇祯期尤其需要奏疏、地方、朝鲜及满文材料交叉。 因此本段仅据材料确认事件的时序、行动和可见后果，并把政策文本与地方实践、报称数字和社会经历分开处理。

\eventanchor{ML-1609-001}{明·1609（万历三十七年）}
\paragraph{1609（万历三十七年）：贾运河竣工} 这项记录把本章所见的制度运作落实到具体年份。其形成背景是：财政征解、交通工程或货币支付的压力使相关措施进入政务。 它带来的、而且仍可由后续材料追踪的改变是：其法定安排不等于地方执行，后续须以地域材料追踪负担与成效。 《神宗实录》万历三十七年条；《明会典》相关条目；地方志、黄册/契约/河工材料 提供相互检验的起点；但译写候选以编年材料定位；实录记载反映朝廷选择，崇祯期须另以奏疏、地方及敌对政权材料交叉。 因此本段仅据材料确认事件的时序、行动和可见后果，并把政策文本与地方实践、报称数字和社会经历分开处理。

\eventanchor{ML-1609-002}{明·1609（万历三十七年）}
\paragraph{1609（万历三十七年）：矿税、加派、漕运和辽饷的本年奏疏与账目提供财政压力线索，预算、欠额和实解不得混同。} 这项记录把本章所见的制度运作落实到具体年份。其形成背景是：财政征解、交通工程或货币支付的压力使相关措施进入政务。 它带来的、而且仍可由后续材料追踪的改变是：其法定安排不等于地方执行，后续须以地域材料追踪负担与成效。 《神宗实录》万历三十七年条；《明会典》相关条目；地方志、黄册/契约/河工材料 提供相互检验的起点；但桥接条目只标示可回查的年度问题与材料链；实录有中央选择性，崇祯期尤其需要奏疏、地方、朝鲜及满文材料交叉。 因此本段仅据材料确认事件的时序、行动和可见后果，并把政策文本与地方实践、报称数字和社会经历分开处理。

\eventanchor{ML-1609-003}{明·1609（万历三十七年）}
\paragraph{1609（万历三十七年）：万历末至天启初的立储、内阁和言路材料标记中枢运作的堵塞与调整，派系归类需回到具体文书。（材料核查重点：诏令或奏议的形成与发受链）} 这项记录把本章所见的制度运作落实到具体年份。其形成背景是：继承、官僚协调或皇权运作的压力使问题进入朝廷程序。 它带来的、而且仍可由后续材料追踪的改变是：它影响后续人事与政策议程；动机和派系标签不能由单一叙事裁定。 《神宗实录》万历三十七年条；《明史·本纪》及相关列传；诏令、奏疏或会典版本 提供相互检验的起点；但桥接条目只标示可回查的年度问题与材料链；实录有中央选择性，崇祯期尤其需要奏疏、地方、朝鲜及满文材料交叉。；本条特别核对诏令或奏议的形成与发受链。 因此本段仅据材料确认事件的时序、行动和可见后果，并把政策文本与地方实践、报称数字和社会经历分开处理。

\subsection{1609年前后的续接}

\eventanchor{ML-1609-004}{明·1609（万历三十七年）}
\paragraph{1609（万历三十七年）：辽东战事、边镇防务和西南兵事的本年军令记录资源调动，战报数字应与朝鲜、满文或地方材料比照。} 这项记录把本章所见的制度运作落实到具体年份。其形成背景是：前期边防、地方武装或跨区域资源调动的压力推动了该行动。 它带来的、而且仍可由后续材料追踪的改变是：它改变了后续调兵、补给或地方秩序的条件；战果和伤亡须分开核对。 《神宗实录》万历三十七年条；《明史·兵志》及相关列传；边镇、地方志或对方材料 提供相互检验的起点；但桥接条目只标示可回查的年度问题与材料链；实录有中央选择性，崇祯期尤其需要奏疏、地方、朝鲜及满文材料交叉。 因此本段仅据材料确认事件的时序、行动和可见后果，并把政策文本与地方实践、报称数字和社会经历分开处理。

\eventanchor{ML-1609-005}{明·1609（万历三十七年）}
\paragraph{1609（万历三十七年）：万历末至天启初的立储、内阁和言路材料标记中枢运作的堵塞与调整，派系归类需回到具体文书。（材料核查重点：报称信息的来源、抄传与行政分类）} 这项记录把本章所见的制度运作落实到具体年份。其形成背景是：继承、官僚协调或皇权运作的压力使问题进入朝廷程序。 它带来的、而且仍可由后续材料追踪的改变是：它影响后续人事与政策议程；动机和派系标签不能由单一叙事裁定。 《神宗实录》万历三十七年条；《明史·本纪》及相关列传；诏令、奏疏或会典版本 提供相互检验的起点；但桥接条目只标示可回查的年度问题与材料链；实录有中央选择性，崇祯期尤其需要奏疏、地方、朝鲜及满文材料交叉。；本条特别核对报称信息的来源、抄传与行政分类。 因此本段仅据材料确认事件的时序、行动和可见后果，并把政策文本与地方实践、报称数字和社会经历分开处理。

\eventanchor{ML-1609-006}{明·1609（万历三十七年）}
\paragraph{1609（万历三十七年）：本年灾害、疫病或地方赈济线索应以受灾地区文书和地方志复核，中央奏报不能覆盖全部社会。} 这项记录把本章所见的制度运作落实到具体年份。其形成背景是：自然风险与地方报告促成朝廷或地方的应对记录。 它带来的、而且仍可由后续材料追踪的改变是：赈济、迁徙和损失的实际范围仍须以受灾地区材料分别核验。 《神宗实录》万历三十七年条；《明史·五行志》；受灾府县地方志与奏疏 提供相互检验的起点；但桥接条目只标示可回查的年度问题与材料链；实录有中央选择性，崇祯期尤其需要奏疏、地方、朝鲜及满文材料交叉。 因此本段仅据材料确认事件的时序、行动和可见后果，并把政策文本与地方实践、报称数字和社会经历分开处理。

\eventanchor{ML-1610-001}{明·1610（万历三十八年）}
\paragraph{1610（万历三十八年）：李约瑟估计，欧洲文明在天文学和物理学方面大约在这个时候超越了中国} 这项记录把本章所见的制度运作落实到具体年份。其形成背景是：继承、官僚协调或皇权运作的压力使问题进入朝廷程序。 它带来的、而且仍可由后续材料追踪的改变是：它影响后续人事与政策议程；动机和派系标签不能由单一叙事裁定。 《神宗实录》万历三十八年条；《明史·本纪》及相关列传；诏令、奏疏或会典版本 提供相互检验的起点；但译写候选以编年材料定位；实录记载反映朝廷选择，崇祯期须另以奏疏、地方及敌对政权材料交叉。 因此本段仅据材料确认事件的时序、行动和可见后果，并把政策文本与地方实践、报称数字和社会经历分开处理。

\eventanchor{ML-1610-002}{明·1610（万历三十八年）}
\paragraph{1610（万历三十八年）：辽东战事、边镇防务和西南兵事的本年军令记录资源调动，战报数字应与朝鲜、满文或地方材料比照。} 这项记录把本章所见的制度运作落实到具体年份。其形成背景是：前期边防、地方武装或跨区域资源调动的压力推动了该行动。 它带来的、而且仍可由后续材料追踪的改变是：它改变了后续调兵、补给或地方秩序的条件；战果和伤亡须分开核对。 《神宗实录》万历三十八年条；《明史·兵志》及相关列传；边镇、地方志或对方材料 提供相互检验的起点；但桥接条目只标示可回查的年度问题与材料链；实录有中央选择性，崇祯期尤其需要奏疏、地方、朝鲜及满文材料交叉。 因此本段仅据材料确认事件的时序、行动和可见后果，并把政策文本与地方实践、报称数字和社会经历分开处理。

\eventanchor{ML-1610-003}{明·1610（万历三十八年）}
\paragraph{1610（万历三十八年）：矿税、加派、漕运和辽饷的本年奏疏与账目提供财政压力线索，预算、欠额和实解不得混同。（材料核查重点：诏令或奏议的形成与发受链）} 这项记录把本章所见的制度运作落实到具体年份。其形成背景是：财政征解、交通工程或货币支付的压力使相关措施进入政务。 它带来的、而且仍可由后续材料追踪的改变是：其法定安排不等于地方执行，后续须以地域材料追踪负担与成效。 《神宗实录》万历三十八年条；《明会典》相关条目；地方志、黄册/契约/河工材料 提供相互检验的起点；但桥接条目只标示可回查的年度问题与材料链；实录有中央选择性，崇祯期尤其需要奏疏、地方、朝鲜及满文材料交叉。；本条特别核对诏令或奏议的形成与发受链。 因此本段仅据材料确认事件的时序、行动和可见后果，并把政策文本与地方实践、报称数字和社会经历分开处理。

\eventanchor{ML-1610-004}{明·1610（万历三十八年）}
\paragraph{1610（万历三十八年）：后金兴起、朝鲜交涉和海上网络的本年多方材料揭示区域关系变化，政权名称与译名须按原文说明。} 这项记录把本章所见的制度运作落实到具体年份。其形成背景是：朝贡名义、边境安全与港口贸易的利益使交涉持续发生。 它带来的、而且仍可由后续材料追踪的改变是：它为后续册封、互市或海防安排留下文书与跨政权比较线索。 《神宗实录》万历三十八年条；《明史·外国传》；《朝鲜王朝实录》、琉球《历代宝案》或海外档案 提供相互检验的起点；但桥接条目只标示可回查的年度问题与材料链；实录有中央选择性，崇祯期尤其需要奏疏、地方、朝鲜及满文材料交叉。 因此本段仅据材料确认事件的时序、行动和可见后果，并把政策文本与地方实践、报称数字和社会经历分开处理。

\eventanchor{ML-1610-005}{明·1610（万历三十八年）}
\paragraph{1610（万历三十八年）：矿税、加派、漕运和辽饷的本年奏疏与账目提供财政压力线索，预算、欠额和实解不得混同。（材料核查重点：报称信息的来源、抄传与行政分类）} 这项记录把本章所见的制度运作落实到具体年份。其形成背景是：财政征解、交通工程或货币支付的压力使相关措施进入政务。 它带来的、而且仍可由后续材料追踪的改变是：其法定安排不等于地方执行，后续须以地域材料追踪负担与成效。 《神宗实录》万历三十八年条；《明会典》相关条目；地方志、黄册/契约/河工材料 提供相互检验的起点；但桥接条目只标示可回查的年度问题与材料链；实录有中央选择性，崇祯期尤其需要奏疏、地方、朝鲜及满文材料交叉。；本条特别核对报称信息的来源、抄传与行政分类。 因此本段仅据材料确认事件的时序、行动和可见后果，并把政策文本与地方实践、报称数字和社会经历分开处理。

\eventanchor{ML-1610-006}{明·1610（万历三十八年）}
\paragraph{1610（万历三十八年）：万历末至天启初的立储、内阁和言路材料标记中枢运作的堵塞与调整，派系归类需回到具体文书。} 这项记录把本章所见的制度运作落实到具体年份。其形成背景是：继承、官僚协调或皇权运作的压力使问题进入朝廷程序。 它带来的、而且仍可由后续材料追踪的改变是：它影响后续人事与政策议程；动机和派系标签不能由单一叙事裁定。 《神宗实录》万历三十八年条；《明史·本纪》及相关列传；诏令、奏疏或会典版本 提供相互检验的起点；但桥接条目只标示可回查的年度问题与材料链；实录有中央选择性，崇祯期尤其需要奏疏、地方、朝鲜及满文材料交叉。 因此本段仅据材料确认事件的时序、行动和可见后果，并把政策文本与地方实践、报称数字和社会经历分开处理。

\eventanchor{ML-1611-001}{明·1611（万历三十九年）}
\paragraph{1611（万历三十九年）：辽东战事、边镇防务和西南兵事的本年军令记录资源调动，战报数字应与朝鲜、满文或地方材料比照。（材料核查重点：诏令或奏议的形成与发受链）} 这项记录把本章所见的制度运作落实到具体年份。其形成背景是：前期边防、地方武装或跨区域资源调动的压力推动了该行动。 它带来的、而且仍可由后续材料追踪的改变是：它改变了后续调兵、补给或地方秩序的条件；战果和伤亡须分开核对。 《神宗实录》万历三十九年条；《明史·兵志》及相关列传；边镇、地方志或对方材料 提供相互检验的起点；但桥接条目只标示可回查的年度问题与材料链；实录有中央选择性，崇祯期尤其需要奏疏、地方、朝鲜及满文材料交叉。；本条特别核对诏令或奏议的形成与发受链。 因此本段仅据材料确认事件的时序、行动和可见后果，并把政策文本与地方实践、报称数字和社会经历分开处理。

\eventanchor{ML-1611-002}{明·1611（万历三十九年）}
\paragraph{1611（万历三十九年）：后金兴起、朝鲜交涉和海上网络的本年多方材料揭示区域关系变化，政权名称与译名须按原文说明。} 这项记录把本章所见的制度运作落实到具体年份。其形成背景是：朝贡名义、边境安全与港口贸易的利益使交涉持续发生。 它带来的、而且仍可由后续材料追踪的改变是：它为后续册封、互市或海防安排留下文书与跨政权比较线索。 《神宗实录》万历三十九年条；《明史·外国传》；《朝鲜王朝实录》、琉球《历代宝案》或海外档案 提供相互检验的起点；但桥接条目只标示可回查的年度问题与材料链；实录有中央选择性，崇祯期尤其需要奏疏、地方、朝鲜及满文材料交叉。 因此本段仅据材料确认事件的时序、行动和可见后果，并把政策文本与地方实践、报称数字和社会经历分开处理。

\eventanchor{ML-1611-003}{明·1611（万历三十九年）}
\paragraph{1611（万历三十九年）：辽东战事、边镇防务和西南兵事的本年军令记录资源调动，战报数字应与朝鲜、满文或地方材料比照。（材料核查重点：报称信息的来源、抄传与行政分类）} 这项记录把本章所见的制度运作落实到具体年份。其形成背景是：前期边防、地方武装或跨区域资源调动的压力推动了该行动。 它带来的、而且仍可由后续材料追踪的改变是：它改变了后续调兵、补给或地方秩序的条件；战果和伤亡须分开核对。 《神宗实录》万历三十九年条；《明史·兵志》及相关列传；边镇、地方志或对方材料 提供相互检验的起点；但桥接条目只标示可回查的年度问题与材料链；实录有中央选择性，崇祯期尤其需要奏疏、地方、朝鲜及满文材料交叉。；本条特别核对报称信息的来源、抄传与行政分类。 因此本段仅据材料确认事件的时序、行动和可见后果，并把政策文本与地方实践、报称数字和社会经历分开处理。

\eventanchor{ML-1611-004}{明·1611（万历三十九年）}
\paragraph{1611（万历三十九年）：东林、书院、刻书和宗教活动的本年材料记录精英网络，社团存在不等于一致政治立场。} 这项记录把本章所见的制度运作落实到具体年份。其形成背景是：制度、教育或知识传播的需要推动了相关行动或文本形成。 它带来的、而且仍可由后续材料追踪的改变是：它提供制度和传播史的锚点，不能据此推断社会整体接受。 《神宗实录》万历三十九年条；《明史·选举志》或《艺文志》；文集、碑刻或书目材料 提供相互检验的起点；但桥接条目只标示可回查的年度问题与材料链；实录有中央选择性，崇祯期尤其需要奏疏、地方、朝鲜及满文材料交叉。 因此本段仅据材料确认事件的时序、行动和可见后果，并把政策文本与地方实践、报称数字和社会经历分开处理。

\eventanchor{ML-1611-005}{明·1611（万历三十九年）}
\paragraph{1611（万历三十九年）：辽东战事、边镇防务和西南兵事的本年军令记录资源调动，战报数字应与朝鲜、满文或地方材料比照。（材料核查重点：同年地方材料所见的执行范围）} 这项记录把本章所见的制度运作落实到具体年份。其形成背景是：前期边防、地方武装或跨区域资源调动的压力推动了该行动。 它带来的、而且仍可由后续材料追踪的改变是：它改变了后续调兵、补给或地方秩序的条件；战果和伤亡须分开核对。 《神宗实录》万历三十九年条；《明史·兵志》及相关列传；边镇、地方志或对方材料 提供相互检验的起点；但桥接条目只标示可回查的年度问题与材料链；实录有中央选择性，崇祯期尤其需要奏疏、地方、朝鲜及满文材料交叉。；本条特别核对同年地方材料所见的执行范围。 因此本段仅据材料确认事件的时序、行动和可见后果，并把政策文本与地方实践、报称数字和社会经历分开处理。

\eventanchor{ML-1611-006}{明·1611（万历三十九年）}
\paragraph{1611（万历三十九年）：矿税、加派、漕运和辽饷的本年奏疏与账目提供财政压力线索，预算、欠额和实解不得混同。} 这项记录把本章所见的制度运作落实到具体年份。其形成背景是：财政征解、交通工程或货币支付的压力使相关措施进入政务。 它带来的、而且仍可由后续材料追踪的改变是：其法定安排不等于地方执行，后续须以地域材料追踪负担与成效。 《神宗实录》万历三十九年条；《明会典》相关条目；地方志、黄册/契约/河工材料 提供相互检验的起点；但桥接条目只标示可回查的年度问题与材料链；实录有中央选择性，崇祯期尤其需要奏疏、地方、朝鲜及满文材料交叉。 因此本段仅据材料确认事件的时序、行动和可见后果，并把政策文本与地方实践、报称数字和社会经历分开处理。

\eventanchor{ML-1612-001}{明·1612（万历四十年）}
\paragraph{1612（万历四十年）：后金兴起、朝鲜交涉和海上网络的本年多方材料揭示区域关系变化，政权名称与译名须按原文说明。（材料核查重点：诏令或奏议的形成与发受链）} 这项记录把本章所见的制度运作落实到具体年份。其形成背景是：朝贡名义、边境安全与港口贸易的利益使交涉持续发生。 它带来的、而且仍可由后续材料追踪的改变是：它为后续册封、互市或海防安排留下文书与跨政权比较线索。 《神宗实录》万历四十年条；《明史·外国传》；《朝鲜王朝实录》、琉球《历代宝案》或海外档案 提供相互检验的起点；但桥接条目只标示可回查的年度问题与材料链；实录有中央选择性，崇祯期尤其需要奏疏、地方、朝鲜及满文材料交叉。；本条特别核对诏令或奏议的形成与发受链。 因此本段仅据材料确认事件的时序、行动和可见后果，并把政策文本与地方实践、报称数字和社会经历分开处理。

\eventanchor{ML-1612-002}{明·1612（万历四十年）}
\paragraph{1612（万历四十年）：东林、书院、刻书和宗教活动的本年材料记录精英网络，社团存在不等于一致政治立场。} 这项记录把本章所见的制度运作落实到具体年份。其形成背景是：制度、教育或知识传播的需要推动了相关行动或文本形成。 它带来的、而且仍可由后续材料追踪的改变是：它提供制度和传播史的锚点，不能据此推断社会整体接受。 《神宗实录》万历四十年条；《明史·选举志》或《艺文志》；文集、碑刻或书目材料 提供相互检验的起点；但桥接条目只标示可回查的年度问题与材料链；实录有中央选择性，崇祯期尤其需要奏疏、地方、朝鲜及满文材料交叉。 因此本段仅据材料确认事件的时序、行动和可见后果，并把政策文本与地方实践、报称数字和社会经历分开处理。

\eventanchor{ML-1612-003}{明·1612（万历四十年）}
\paragraph{1612（万历四十年）：后金兴起、朝鲜交涉和海上网络的本年多方材料揭示区域关系变化，政权名称与译名须按原文说明。（材料核查重点：报称信息的来源、抄传与行政分类）} 这项记录把本章所见的制度运作落实到具体年份。其形成背景是：朝贡名义、边境安全与港口贸易的利益使交涉持续发生。 它带来的、而且仍可由后续材料追踪的改变是：它为后续册封、互市或海防安排留下文书与跨政权比较线索。 《神宗实录》万历四十年条；《明史·外国传》；《朝鲜王朝实录》、琉球《历代宝案》或海外档案 提供相互检验的起点；但桥接条目只标示可回查的年度问题与材料链；实录有中央选择性，崇祯期尤其需要奏疏、地方、朝鲜及满文材料交叉。；本条特别核对报称信息的来源、抄传与行政分类。 因此本段仅据材料确认事件的时序、行动和可见后果，并把政策文本与地方实践、报称数字和社会经历分开处理。

\eventanchor{ML-1612-004}{明·1612（万历四十年）}
\paragraph{1612（万历四十年）：本年灾害、疫病或地方赈济线索应以受灾地区文书和地方志复核，中央奏报不能覆盖全部社会。} 这项记录把本章所见的制度运作落实到具体年份。其形成背景是：自然风险与地方报告促成朝廷或地方的应对记录。 它带来的、而且仍可由后续材料追踪的改变是：赈济、迁徙和损失的实际范围仍须以受灾地区材料分别核验。 《神宗实录》万历四十年条；《明史·五行志》；受灾府县地方志与奏疏 提供相互检验的起点；但桥接条目只标示可回查的年度问题与材料链；实录有中央选择性，崇祯期尤其需要奏疏、地方、朝鲜及满文材料交叉。 因此本段仅据材料确认事件的时序、行动和可见后果，并把政策文本与地方实践、报称数字和社会经历分开处理。

\eventanchor{ML-1612-005}{明·1612（万历四十年）}
\paragraph{1612（万历四十年）：后金兴起、朝鲜交涉和海上网络的本年多方材料揭示区域关系变化，政权名称与译名须按原文说明。（材料核查重点：同年地方材料所见的执行范围）} 这项记录把本章所见的制度运作落实到具体年份。其形成背景是：朝贡名义、边境安全与港口贸易的利益使交涉持续发生。 它带来的、而且仍可由后续材料追踪的改变是：它为后续册封、互市或海防安排留下文书与跨政权比较线索。 《神宗实录》万历四十年条；《明史·外国传》；《朝鲜王朝实录》、琉球《历代宝案》或海外档案 提供相互检验的起点；但桥接条目只标示可回查的年度问题与材料链；实录有中央选择性，崇祯期尤其需要奏疏、地方、朝鲜及满文材料交叉。；本条特别核对同年地方材料所见的执行范围。 因此本段仅据材料确认事件的时序、行动和可见后果，并把政策文本与地方实践、报称数字和社会经历分开处理。

\subsection{1612年前后的续接}

\eventanchor{ML-1612-006}{明·1612（万历四十年）}
\paragraph{1612（万历四十年）：辽东战事、边镇防务和西南兵事的本年军令记录资源调动，战报数字应与朝鲜、满文或地方材料比照。} 这项记录把本章所见的制度运作落实到具体年份。其形成背景是：前期边防、地方武装或跨区域资源调动的压力推动了该行动。 它带来的、而且仍可由后续材料追踪的改变是：它改变了后续调兵、补给或地方秩序的条件；战果和伤亡须分开核对。 《神宗实录》万历四十年条；《明史·兵志》及相关列传；边镇、地方志或对方材料 提供相互检验的起点；但桥接条目只标示可回查的年度问题与材料链；实录有中央选择性，崇祯期尤其需要奏疏、地方、朝鲜及满文材料交叉。 因此本段仅据材料确认事件的时序、行动和可见后果，并把政策文本与地方实践、报称数字和社会经历分开处理。

\eventanchor{ML-1613-001}{明·1613（万历四十一年）}
\paragraph{1613（万历四十一年）：东林、书院、刻书和宗教活动的本年材料记录精英网络，社团存在不等于一致政治立场。（材料核查重点：诏令或奏议的形成与发受链）} 这项记录把本章所见的制度运作落实到具体年份。其形成背景是：制度、教育或知识传播的需要推动了相关行动或文本形成。 它带来的、而且仍可由后续材料追踪的改变是：它提供制度和传播史的锚点，不能据此推断社会整体接受。 《神宗实录》万历四十一年条；《明史·选举志》或《艺文志》；文集、碑刻或书目材料 提供相互检验的起点；但桥接条目只标示可回查的年度问题与材料链；实录有中央选择性，崇祯期尤其需要奏疏、地方、朝鲜及满文材料交叉。；本条特别核对诏令或奏议的形成与发受链。 因此本段仅据材料确认事件的时序、行动和可见后果，并把政策文本与地方实践、报称数字和社会经历分开处理。

\eventanchor{ML-1613-002}{明·1613（万历四十一年）}
\paragraph{1613（万历四十一年）：本年灾害、疫病或地方赈济线索应以受灾地区文书和地方志复核，中央奏报不能覆盖全部社会。} 这项记录把本章所见的制度运作落实到具体年份。其形成背景是：自然风险与地方报告促成朝廷或地方的应对记录。 它带来的、而且仍可由后续材料追踪的改变是：赈济、迁徙和损失的实际范围仍须以受灾地区材料分别核验。 《神宗实录》万历四十一年条；《明史·五行志》；受灾府县地方志与奏疏 提供相互检验的起点；但桥接条目只标示可回查的年度问题与材料链；实录有中央选择性，崇祯期尤其需要奏疏、地方、朝鲜及满文材料交叉。 因此本段仅据材料确认事件的时序、行动和可见后果，并把政策文本与地方实践、报称数字和社会经历分开处理。

\eventanchor{ML-1613-003}{明·1613（万历四十一年）}
\paragraph{1613（万历四十一年）：东林、书院、刻书和宗教活动的本年材料记录精英网络，社团存在不等于一致政治立场。（材料核查重点：报称信息的来源、抄传与行政分类）} 这项记录把本章所见的制度运作落实到具体年份。其形成背景是：制度、教育或知识传播的需要推动了相关行动或文本形成。 它带来的、而且仍可由后续材料追踪的改变是：它提供制度和传播史的锚点，不能据此推断社会整体接受。 《神宗实录》万历四十一年条；《明史·选举志》或《艺文志》；文集、碑刻或书目材料 提供相互检验的起点；但桥接条目只标示可回查的年度问题与材料链；实录有中央选择性，崇祯期尤其需要奏疏、地方、朝鲜及满文材料交叉。；本条特别核对报称信息的来源、抄传与行政分类。 因此本段仅据材料确认事件的时序、行动和可见后果，并把政策文本与地方实践、报称数字和社会经历分开处理。

\eventanchor{ML-1613-004}{明·1613（万历四十一年）}
\paragraph{1613（万历四十一年）：万历末至天启初的立储、内阁和言路材料标记中枢运作的堵塞与调整，派系归类需回到具体文书。} 这项记录把本章所见的制度运作落实到具体年份。其形成背景是：继承、官僚协调或皇权运作的压力使问题进入朝廷程序。 它带来的、而且仍可由后续材料追踪的改变是：它影响后续人事与政策议程；动机和派系标签不能由单一叙事裁定。 《神宗实录》万历四十一年条；《明史·本纪》及相关列传；诏令、奏疏或会典版本 提供相互检验的起点；但桥接条目只标示可回查的年度问题与材料链；实录有中央选择性，崇祯期尤其需要奏疏、地方、朝鲜及满文材料交叉。 因此本段仅据材料确认事件的时序、行动和可见后果，并把政策文本与地方实践、报称数字和社会经历分开处理。

\eventanchor{ML-1613-005}{明·1613（万历四十一年）}
\paragraph{1613（万历四十一年）：东林、书院、刻书和宗教活动的本年材料记录精英网络，社团存在不等于一致政治立场。（材料核查重点：同年地方材料所见的执行范围）} 这项记录把本章所见的制度运作落实到具体年份。其形成背景是：制度、教育或知识传播的需要推动了相关行动或文本形成。 它带来的、而且仍可由后续材料追踪的改变是：它提供制度和传播史的锚点，不能据此推断社会整体接受。 《神宗实录》万历四十一年条；《明史·选举志》或《艺文志》；文集、碑刻或书目材料 提供相互检验的起点；但桥接条目只标示可回查的年度问题与材料链；实录有中央选择性，崇祯期尤其需要奏疏、地方、朝鲜及满文材料交叉。；本条特别核对同年地方材料所见的执行范围。 因此本段仅据材料确认事件的时序、行动和可见后果，并把政策文本与地方实践、报称数字和社会经历分开处理。

\eventanchor{ML-1613-006}{明·1613（万历四十一年）}
\paragraph{1613（万历四十一年）：后金兴起、朝鲜交涉和海上网络的本年多方材料揭示区域关系变化，政权名称与译名须按原文说明。} 这项记录把本章所见的制度运作落实到具体年份。其形成背景是：朝贡名义、边境安全与港口贸易的利益使交涉持续发生。 它带来的、而且仍可由后续材料追踪的改变是：它为后续册封、互市或海防安排留下文书与跨政权比较线索。 《神宗实录》万历四十一年条；《明史·外国传》；《朝鲜王朝实录》、琉球《历代宝案》或海外档案 提供相互检验的起点；但桥接条目只标示可回查的年度问题与材料链；实录有中央选择性，崇祯期尤其需要奏疏、地方、朝鲜及满文材料交叉。 因此本段仅据材料确认事件的时序、行动和可见后果，并把政策文本与地方实践、报称数字和社会经历分开处理。

\eventanchor{ML-1614-001}{明·1614（万历四十二年）}
\paragraph{1614（万历四十二年）：本年灾害、疫病或地方赈济线索应以受灾地区文书和地方志复核，中央奏报不能覆盖全部社会。（材料核查重点：诏令或奏议的形成与发受链）} 这项记录把本章所见的制度运作落实到具体年份。其形成背景是：自然风险与地方报告促成朝廷或地方的应对记录。 它带来的、而且仍可由后续材料追踪的改变是：赈济、迁徙和损失的实际范围仍须以受灾地区材料分别核验。 《神宗实录》万历四十二年条；《明史·五行志》；受灾府县地方志与奏疏 提供相互检验的起点；但桥接条目只标示可回查的年度问题与材料链；实录有中央选择性，崇祯期尤其需要奏疏、地方、朝鲜及满文材料交叉。；本条特别核对诏令或奏议的形成与发受链。 因此本段仅据材料确认事件的时序、行动和可见后果，并把政策文本与地方实践、报称数字和社会经历分开处理。

\eventanchor{ML-1614-002}{明·1614（万历四十二年）}
\paragraph{1614（万历四十二年）：万历末至天启初的立储、内阁和言路材料标记中枢运作的堵塞与调整，派系归类需回到具体文书。} 这项记录把本章所见的制度运作落实到具体年份。其形成背景是：继承、官僚协调或皇权运作的压力使问题进入朝廷程序。 它带来的、而且仍可由后续材料追踪的改变是：它影响后续人事与政策议程；动机和派系标签不能由单一叙事裁定。 《神宗实录》万历四十二年条；《明史·本纪》及相关列传；诏令、奏疏或会典版本 提供相互检验的起点；但桥接条目只标示可回查的年度问题与材料链；实录有中央选择性，崇祯期尤其需要奏疏、地方、朝鲜及满文材料交叉。 因此本段仅据材料确认事件的时序、行动和可见后果，并把政策文本与地方实践、报称数字和社会经历分开处理。

\eventanchor{ML-1614-003}{明·1614（万历四十二年）}
\paragraph{1614（万历四十二年）：本年灾害、疫病或地方赈济线索应以受灾地区文书和地方志复核，中央奏报不能覆盖全部社会。（材料核查重点：报称信息的来源、抄传与行政分类）} 这项记录把本章所见的制度运作落实到具体年份。其形成背景是：自然风险与地方报告促成朝廷或地方的应对记录。 它带来的、而且仍可由后续材料追踪的改变是：赈济、迁徙和损失的实际范围仍须以受灾地区材料分别核验。 《神宗实录》万历四十二年条；《明史·五行志》；受灾府县地方志与奏疏 提供相互检验的起点；但桥接条目只标示可回查的年度问题与材料链；实录有中央选择性，崇祯期尤其需要奏疏、地方、朝鲜及满文材料交叉。；本条特别核对报称信息的来源、抄传与行政分类。 因此本段仅据材料确认事件的时序、行动和可见后果，并把政策文本与地方实践、报称数字和社会经历分开处理。

\eventanchor{ML-1614-004}{明·1614（万历四十二年）}
\paragraph{1614（万历四十二年）：矿税、加派、漕运和辽饷的本年奏疏与账目提供财政压力线索，预算、欠额和实解不得混同。} 这项记录把本章所见的制度运作落实到具体年份。其形成背景是：财政征解、交通工程或货币支付的压力使相关措施进入政务。 它带来的、而且仍可由后续材料追踪的改变是：其法定安排不等于地方执行，后续须以地域材料追踪负担与成效。 《神宗实录》万历四十二年条；《明会典》相关条目；地方志、黄册/契约/河工材料 提供相互检验的起点；但桥接条目只标示可回查的年度问题与材料链；实录有中央选择性，崇祯期尤其需要奏疏、地方、朝鲜及满文材料交叉。 因此本段仅据材料确认事件的时序、行动和可见后果，并把政策文本与地方实践、报称数字和社会经历分开处理。

\eventanchor{ML-1614-005}{明·1614（万历四十二年）}
\paragraph{1614（万历四十二年）：本年灾害、疫病或地方赈济线索应以受灾地区文书和地方志复核，中央奏报不能覆盖全部社会。（材料核查重点：同年地方材料所见的执行范围）} 这项记录把本章所见的制度运作落实到具体年份。其形成背景是：自然风险与地方报告促成朝廷或地方的应对记录。 它带来的、而且仍可由后续材料追踪的改变是：赈济、迁徙和损失的实际范围仍须以受灾地区材料分别核验。 《神宗实录》万历四十二年条；《明史·五行志》；受灾府县地方志与奏疏 提供相互检验的起点；但桥接条目只标示可回查的年度问题与材料链；实录有中央选择性，崇祯期尤其需要奏疏、地方、朝鲜及满文材料交叉。；本条特别核对同年地方材料所见的执行范围。 因此本段仅据材料确认事件的时序、行动和可见后果，并把政策文本与地方实践、报称数字和社会经历分开处理。

\eventanchor{ML-1614-006}{明·1614（万历四十二年）}
\paragraph{1614（万历四十二年）：东林、书院、刻书和宗教活动的本年材料记录精英网络，社团存在不等于一致政治立场。} 这项记录把本章所见的制度运作落实到具体年份。其形成背景是：制度、教育或知识传播的需要推动了相关行动或文本形成。 它带来的、而且仍可由后续材料追踪的改变是：它提供制度和传播史的锚点，不能据此推断社会整体接受。 《神宗实录》万历四十二年条；《明史·选举志》或《艺文志》；文集、碑刻或书目材料 提供相互检验的起点；但桥接条目只标示可回查的年度问题与材料链；实录有中央选择性，崇祯期尤其需要奏疏、地方、朝鲜及满文材料交叉。 因此本段仅据材料确认事件的时序、行动和可见后果，并把政策文本与地方实践、报称数字和社会经历分开处理。

\eventanchor{ML-1615-001}{明·1615（万历四十三年）}
\paragraph{1615（万历四十三年）：努尔哈赤派遣最后一位朝贡使者前往北京} 这项记录把本章所见的制度运作落实到具体年份。其形成背景是：继承、官僚协调或皇权运作的压力使问题进入朝廷程序。 它带来的、而且仍可由后续材料追踪的改变是：它影响后续人事与政策议程；动机和派系标签不能由单一叙事裁定。 《神宗实录》万历四十三年条；《明史·本纪》及相关列传；诏令、奏疏或会典版本 提供相互检验的起点；但译写候选以编年材料定位；实录记载反映朝廷选择，崇祯期须另以奏疏、地方及敌对政权材料交叉。 因此本段仅据材料确认事件的时序、行动和可见后果，并把政策文本与地方实践、报称数字和社会经历分开处理。

\eventanchor{ML-1615-002}{明·1615（万历四十三年）}
\paragraph{1615（万历四十三年）：矿税、加派、漕运和辽饷的本年奏疏与账目提供财政压力线索，预算、欠额和实解不得混同。} 这项记录把本章所见的制度运作落实到具体年份。其形成背景是：财政征解、交通工程或货币支付的压力使相关措施进入政务。 它带来的、而且仍可由后续材料追踪的改变是：其法定安排不等于地方执行，后续须以地域材料追踪负担与成效。 《神宗实录》万历四十三年条；《明会典》相关条目；地方志、黄册/契约/河工材料 提供相互检验的起点；但桥接条目只标示可回查的年度问题与材料链；实录有中央选择性，崇祯期尤其需要奏疏、地方、朝鲜及满文材料交叉。 因此本段仅据材料确认事件的时序、行动和可见后果，并把政策文本与地方实践、报称数字和社会经历分开处理。

\eventanchor{ML-1615-003}{明·1615（万历四十三年）}
\paragraph{1615（万历四十三年）：万历末至天启初的立储、内阁和言路材料标记中枢运作的堵塞与调整，派系归类需回到具体文书。（材料核查重点：诏令或奏议的形成与发受链）} 这项记录把本章所见的制度运作落实到具体年份。其形成背景是：继承、官僚协调或皇权运作的压力使问题进入朝廷程序。 它带来的、而且仍可由后续材料追踪的改变是：它影响后续人事与政策议程；动机和派系标签不能由单一叙事裁定。 《神宗实录》万历四十三年条；《明史·本纪》及相关列传；诏令、奏疏或会典版本 提供相互检验的起点；但桥接条目只标示可回查的年度问题与材料链；实录有中央选择性，崇祯期尤其需要奏疏、地方、朝鲜及满文材料交叉。；本条特别核对诏令或奏议的形成与发受链。 因此本段仅据材料确认事件的时序、行动和可见后果，并把政策文本与地方实践、报称数字和社会经历分开处理。

\eventanchor{ML-1615-004}{明·1615（万历四十三年）}
\paragraph{1615（万历四十三年）：辽东战事、边镇防务和西南兵事的本年军令记录资源调动，战报数字应与朝鲜、满文或地方材料比照。} 这项记录把本章所见的制度运作落实到具体年份。其形成背景是：前期边防、地方武装或跨区域资源调动的压力推动了该行动。 它带来的、而且仍可由后续材料追踪的改变是：它改变了后续调兵、补给或地方秩序的条件；战果和伤亡须分开核对。 《神宗实录》万历四十三年条；《明史·兵志》及相关列传；边镇、地方志或对方材料 提供相互检验的起点；但桥接条目只标示可回查的年度问题与材料链；实录有中央选择性，崇祯期尤其需要奏疏、地方、朝鲜及满文材料交叉。 因此本段仅据材料确认事件的时序、行动和可见后果，并把政策文本与地方实践、报称数字和社会经历分开处理。

\eventanchor{ML-1615-005}{明·1615（万历四十三年）}
\paragraph{1615（万历四十三年）：万历末至天启初的立储、内阁和言路材料标记中枢运作的堵塞与调整，派系归类需回到具体文书。（材料核查重点：报称信息的来源、抄传与行政分类）} 这项记录把本章所见的制度运作落实到具体年份。其形成背景是：继承、官僚协调或皇权运作的压力使问题进入朝廷程序。 它带来的、而且仍可由后续材料追踪的改变是：它影响后续人事与政策议程；动机和派系标签不能由单一叙事裁定。 《神宗实录》万历四十三年条；《明史·本纪》及相关列传；诏令、奏疏或会典版本 提供相互检验的起点；但桥接条目只标示可回查的年度问题与材料链；实录有中央选择性，崇祯期尤其需要奏疏、地方、朝鲜及满文材料交叉。；本条特别核对报称信息的来源、抄传与行政分类。 因此本段仅据材料确认事件的时序、行动和可见后果，并把政策文本与地方实践、报称数字和社会经历分开处理。

\eventanchor{ML-1615-006}{明·1615（万历四十三年）}
\paragraph{1615（万历四十三年）：本年灾害、疫病或地方赈济线索应以受灾地区文书和地方志复核，中央奏报不能覆盖全部社会。} 这项记录把本章所见的制度运作落实到具体年份。其形成背景是：自然风险与地方报告促成朝廷或地方的应对记录。 它带来的、而且仍可由后续材料追踪的改变是：赈济、迁徙和损失的实际范围仍须以受灾地区材料分别核验。 《神宗实录》万历四十三年条；《明史·五行志》；受灾府县地方志与奏疏 提供相互检验的起点；但桥接条目只标示可回查的年度问题与材料链；实录有中央选择性，崇祯期尤其需要奏疏、地方、朝鲜及满文材料交叉。 因此本段仅据材料确认事件的时序、行动和可见后果，并把政策文本与地方实践、报称数字和社会经历分开处理。

\eventanchor{ML-1616-001}{明·1616（万历四十四年）}
\paragraph{1616（万历四十四年）：努尔哈赤宣告后金，又称阿玛嘎爱新古伦} 这项记录把本章所见的制度运作落实到具体年份。其形成背景是：继承、官僚协调或皇权运作的压力使问题进入朝廷程序。 它带来的、而且仍可由后续材料追踪的改变是：它影响后续人事与政策议程；动机和派系标签不能由单一叙事裁定。 《神宗实录》万历四十四年条；《明史·本纪》及相关列传；诏令、奏疏或会典版本 提供相互检验的起点；但译写候选以编年材料定位；实录记载反映朝廷选择，崇祯期须另以奏疏、地方及敌对政权材料交叉。 因此本段仅据材料确认事件的时序、行动和可见后果，并把政策文本与地方实践、报称数字和社会经历分开处理。

\subsection{1616年前后的续接}

\eventanchor{ML-1616-002}{明·1616（万历四十四年）}
\paragraph{1616（万历四十四年）：辽东战事、边镇防务和西南兵事的本年军令记录资源调动，战报数字应与朝鲜、满文或地方材料比照。} 这项记录把本章所见的制度运作落实到具体年份。其形成背景是：前期边防、地方武装或跨区域资源调动的压力推动了该行动。 它带来的、而且仍可由后续材料追踪的改变是：它改变了后续调兵、补给或地方秩序的条件；战果和伤亡须分开核对。 《神宗实录》万历四十四年条；《明史·兵志》及相关列传；边镇、地方志或对方材料 提供相互检验的起点；但桥接条目只标示可回查的年度问题与材料链；实录有中央选择性，崇祯期尤其需要奏疏、地方、朝鲜及满文材料交叉。 因此本段仅据材料确认事件的时序、行动和可见后果，并把政策文本与地方实践、报称数字和社会经历分开处理。

\eventanchor{ML-1616-003}{明·1616（万历四十四年）}
\paragraph{1616（万历四十四年）：矿税、加派、漕运和辽饷的本年奏疏与账目提供财政压力线索，预算、欠额和实解不得混同。（材料核查重点：诏令或奏议的形成与发受链）} 这项记录把本章所见的制度运作落实到具体年份。其形成背景是：财政征解、交通工程或货币支付的压力使相关措施进入政务。 它带来的、而且仍可由后续材料追踪的改变是：其法定安排不等于地方执行，后续须以地域材料追踪负担与成效。 《神宗实录》万历四十四年条；《明会典》相关条目；地方志、黄册/契约/河工材料 提供相互检验的起点；但桥接条目只标示可回查的年度问题与材料链；实录有中央选择性，崇祯期尤其需要奏疏、地方、朝鲜及满文材料交叉。；本条特别核对诏令或奏议的形成与发受链。 因此本段仅据材料确认事件的时序、行动和可见后果，并把政策文本与地方实践、报称数字和社会经历分开处理。

\eventanchor{ML-1616-004}{明·1616（万历四十四年）}
\paragraph{1616（万历四十四年）：后金兴起、朝鲜交涉和海上网络的本年多方材料揭示区域关系变化，政权名称与译名须按原文说明。} 这项记录把本章所见的制度运作落实到具体年份。其形成背景是：朝贡名义、边境安全与港口贸易的利益使交涉持续发生。 它带来的、而且仍可由后续材料追踪的改变是：它为后续册封、互市或海防安排留下文书与跨政权比较线索。 《神宗实录》万历四十四年条；《明史·外国传》；《朝鲜王朝实录》、琉球《历代宝案》或海外档案 提供相互检验的起点；但桥接条目只标示可回查的年度问题与材料链；实录有中央选择性，崇祯期尤其需要奏疏、地方、朝鲜及满文材料交叉。 因此本段仅据材料确认事件的时序、行动和可见后果，并把政策文本与地方实践、报称数字和社会经历分开处理。

\eventanchor{ML-1616-005}{明·1616（万历四十四年）}
\paragraph{1616（万历四十四年）：矿税、加派、漕运和辽饷的本年奏疏与账目提供财政压力线索，预算、欠额和实解不得混同。（材料核查重点：报称信息的来源、抄传与行政分类）} 这项记录把本章所见的制度运作落实到具体年份。其形成背景是：财政征解、交通工程或货币支付的压力使相关措施进入政务。 它带来的、而且仍可由后续材料追踪的改变是：其法定安排不等于地方执行，后续须以地域材料追踪负担与成效。 《神宗实录》万历四十四年条；《明会典》相关条目；地方志、黄册/契约/河工材料 提供相互检验的起点；但桥接条目只标示可回查的年度问题与材料链；实录有中央选择性，崇祯期尤其需要奏疏、地方、朝鲜及满文材料交叉。；本条特别核对报称信息的来源、抄传与行政分类。 因此本段仅据材料确认事件的时序、行动和可见后果，并把政策文本与地方实践、报称数字和社会经历分开处理。

\eventanchor{ML-1616-006}{明·1616（万历四十四年）}
\paragraph{1616（万历四十四年）：万历末至天启初的立储、内阁和言路材料标记中枢运作的堵塞与调整，派系归类需回到具体文书。} 这项记录把本章所见的制度运作落实到具体年份。其形成背景是：继承、官僚协调或皇权运作的压力使问题进入朝廷程序。 它带来的、而且仍可由后续材料追踪的改变是：它影响后续人事与政策议程；动机和派系标签不能由单一叙事裁定。 《神宗实录》万历四十四年条；《明史·本纪》及相关列传；诏令、奏疏或会典版本 提供相互检验的起点；但桥接条目只标示可回查的年度问题与材料链；实录有中央选择性，崇祯期尤其需要奏疏、地方、朝鲜及满文材料交叉。 因此本段仅据材料确认事件的时序、行动和可见后果，并把政策文本与地方实践、报称数字和社会经历分开处理。

\eventanchor{ML-1617-001}{明·1617（万历四十五年）}
\paragraph{1617（万历四十五年）：万历末至天启初的立储、内阁和言路材料标记中枢运作的堵塞与调整，派系归类需回到具体文书。} 这项记录把本章所见的制度运作落实到具体年份。其形成背景是：继承、官僚协调或皇权运作的压力使问题进入朝廷程序。 它带来的、而且仍可由后续材料追踪的改变是：它影响后续人事与政策议程；动机和派系标签不能由单一叙事裁定。 《神宗实录》万历四十五年条；《明史·本纪》及相关列传；诏令、奏疏或会典版本 提供相互检验的起点；但桥接条目只标示可回查的年度问题与材料链；实录有中央选择性，崇祯期尤其需要奏疏、地方、朝鲜及满文材料交叉。 因此本段仅据材料确认事件的时序、行动和可见后果，并把政策文本与地方实践、报称数字和社会经历分开处理。

\eventanchor{ML-1617-002}{明·1617（万历四十五年）}
\paragraph{1617（万历四十五年）：后金兴起、朝鲜交涉和海上网络的本年多方材料揭示区域关系变化，政权名称与译名须按原文说明。} 这项记录把本章所见的制度运作落实到具体年份。其形成背景是：朝贡名义、边境安全与港口贸易的利益使交涉持续发生。 它带来的、而且仍可由后续材料追踪的改变是：它为后续册封、互市或海防安排留下文书与跨政权比较线索。 《神宗实录》万历四十五年条；《明史·外国传》；《朝鲜王朝实录》、琉球《历代宝案》或海外档案 提供相互检验的起点；但桥接条目只标示可回查的年度问题与材料链；实录有中央选择性，崇祯期尤其需要奏疏、地方、朝鲜及满文材料交叉。 因此本段仅据材料确认事件的时序、行动和可见后果，并把政策文本与地方实践、报称数字和社会经历分开处理。

\eventanchor{ML-1617-003}{明·1617（万历四十五年）}
\paragraph{1617（万历四十五年）：辽东战事、边镇防务和西南兵事的本年军令记录资源调动，战报数字应与朝鲜、满文或地方材料比照。（材料核查重点：诏令或奏议的形成与发受链）} 这项记录把本章所见的制度运作落实到具体年份。其形成背景是：前期边防、地方武装或跨区域资源调动的压力推动了该行动。 它带来的、而且仍可由后续材料追踪的改变是：它改变了后续调兵、补给或地方秩序的条件；战果和伤亡须分开核对。 《神宗实录》万历四十五年条；《明史·兵志》及相关列传；边镇、地方志或对方材料 提供相互检验的起点；但桥接条目只标示可回查的年度问题与材料链；实录有中央选择性，崇祯期尤其需要奏疏、地方、朝鲜及满文材料交叉。；本条特别核对诏令或奏议的形成与发受链。 因此本段仅据材料确认事件的时序、行动和可见后果，并把政策文本与地方实践、报称数字和社会经历分开处理。

\eventanchor{ML-1617-004}{明·1617（万历四十五年）}
\paragraph{1617（万历四十五年）：东林、书院、刻书和宗教活动的本年材料记录精英网络，社团存在不等于一致政治立场。} 这项记录把本章所见的制度运作落实到具体年份。其形成背景是：制度、教育或知识传播的需要推动了相关行动或文本形成。 它带来的、而且仍可由后续材料追踪的改变是：它提供制度和传播史的锚点，不能据此推断社会整体接受。 《神宗实录》万历四十五年条；《明史·选举志》或《艺文志》；文集、碑刻或书目材料 提供相互检验的起点；但桥接条目只标示可回查的年度问题与材料链；实录有中央选择性，崇祯期尤其需要奏疏、地方、朝鲜及满文材料交叉。 因此本段仅据材料确认事件的时序、行动和可见后果，并把政策文本与地方实践、报称数字和社会经历分开处理。

\eventanchor{ML-1617-005}{明·1617（万历四十五年）}
\paragraph{1617（万历四十五年）：辽东战事、边镇防务和西南兵事的本年军令记录资源调动，战报数字应与朝鲜、满文或地方材料比照。（材料核查重点：报称信息的来源、抄传与行政分类）} 这项记录把本章所见的制度运作落实到具体年份。其形成背景是：前期边防、地方武装或跨区域资源调动的压力推动了该行动。 它带来的、而且仍可由后续材料追踪的改变是：它改变了后续调兵、补给或地方秩序的条件；战果和伤亡须分开核对。 《神宗实录》万历四十五年条；《明史·兵志》及相关列传；边镇、地方志或对方材料 提供相互检验的起点；但桥接条目只标示可回查的年度问题与材料链；实录有中央选择性，崇祯期尤其需要奏疏、地方、朝鲜及满文材料交叉。；本条特别核对报称信息的来源、抄传与行政分类。 因此本段仅据材料确认事件的时序、行动和可见后果，并把政策文本与地方实践、报称数字和社会经历分开处理。

\eventanchor{ML-1617-006}{明·1617（万历四十五年）}
\paragraph{1617（万历四十五年）：矿税、加派、漕运和辽饷的本年奏疏与账目提供财政压力线索，预算、欠额和实解不得混同。} 这项记录把本章所见的制度运作落实到具体年份。其形成背景是：财政征解、交通工程或货币支付的压力使相关措施进入政务。 它带来的、而且仍可由后续材料追踪的改变是：其法定安排不等于地方执行，后续须以地域材料追踪负担与成效。 《神宗实录》万历四十五年条；《明会典》相关条目；地方志、黄册/契约/河工材料 提供相互检验的起点；但桥接条目只标示可回查的年度问题与材料链；实录有中央选择性，崇祯期尤其需要奏疏、地方、朝鲜及满文材料交叉。 因此本段仅据材料确认事件的时序、行动和可见后果，并把政策文本与地方实践、报称数字和社会经历分开处理。

\eventanchor{ML-1618-001}{明·1618（万历四十六年）五月9日}
\paragraph{1618（万历四十六年）五月9日：抚顺之战：后金占领抚顺} 这项记录把本章所见的制度运作落实到具体年份。其形成背景是：前期边防、地方武装或跨区域资源调动的压力推动了该行动。 它带来的、而且仍可由后续材料追踪的改变是：它改变了后续调兵、补给或地方秩序的条件；战果和伤亡须分开核对。 《神宗实录》万历四十六年条；《明史·兵志》及相关列传；边镇、地方志或对方材料 提供相互检验的起点；但译写候选以编年材料定位；实录记载反映朝廷选择，崇祯期须另以奏疏、地方及敌对政权材料交叉。 因此本段仅据材料确认事件的时序、行动和可见后果，并把政策文本与地方实践、报称数字和社会经历分开处理。

\eventanchor{ML-1618-002}{明·1618（万历四十六年）夏}
\paragraph{1618（万历四十六年）夏：清河之战：后金攻克清河} 这项记录把本章所见的制度运作落实到具体年份。其形成背景是：前期边防、地方武装或跨区域资源调动的压力推动了该行动。 它带来的、而且仍可由后续材料追踪的改变是：它改变了后续调兵、补给或地方秩序的条件；战果和伤亡须分开核对。 《神宗实录》万历四十六年条；《明史·兵志》及相关列传；边镇、地方志或对方材料 提供相互检验的起点；但译写候选以编年材料定位；实录记载反映朝廷选择，崇祯期须另以奏疏、地方及敌对政权材料交叉。 因此本段仅据材料确认事件的时序、行动和可见后果，并把政策文本与地方实践、报称数字和社会经历分开处理。

\eventanchor{ML-1618-003}{明·1618（万历四十六年）}
\paragraph{1618（万历四十六年）：后金兴起、朝鲜交涉和海上网络的本年多方材料揭示区域关系变化，政权名称与译名须按原文说明。（材料核查重点：诏令或奏议的形成与发受链）} 这项记录把本章所见的制度运作落实到具体年份。其形成背景是：朝贡名义、边境安全与港口贸易的利益使交涉持续发生。 它带来的、而且仍可由后续材料追踪的改变是：它为后续册封、互市或海防安排留下文书与跨政权比较线索。 《神宗实录》万历四十六年条；《明史·外国传》；《朝鲜王朝实录》、琉球《历代宝案》或海外档案 提供相互检验的起点；但桥接条目只标示可回查的年度问题与材料链；实录有中央选择性，崇祯期尤其需要奏疏、地方、朝鲜及满文材料交叉。；本条特别核对诏令或奏议的形成与发受链。 因此本段仅据材料确认事件的时序、行动和可见后果，并把政策文本与地方实践、报称数字和社会经历分开处理。

\eventanchor{ML-1618-004}{明·1618（万历四十六年）}
\paragraph{1618（万历四十六年）：本年灾害、疫病或地方赈济线索应以受灾地区文书和地方志复核，中央奏报不能覆盖全部社会。} 这项记录把本章所见的制度运作落实到具体年份。其形成背景是：自然风险与地方报告促成朝廷或地方的应对记录。 它带来的、而且仍可由后续材料追踪的改变是：赈济、迁徙和损失的实际范围仍须以受灾地区材料分别核验。 《神宗实录》万历四十六年条；《明史·五行志》；受灾府县地方志与奏疏 提供相互检验的起点；但桥接条目只标示可回查的年度问题与材料链；实录有中央选择性，崇祯期尤其需要奏疏、地方、朝鲜及满文材料交叉。 因此本段仅据材料确认事件的时序、行动和可见后果，并把政策文本与地方实践、报称数字和社会经历分开处理。

\eventanchor{ML-1618-005}{明·1618（万历四十六年）}
\paragraph{1618（万历四十六年）：后金兴起、朝鲜交涉和海上网络的本年多方材料揭示区域关系变化，政权名称与译名须按原文说明。（材料核查重点：报称信息的来源、抄传与行政分类）} 这项记录把本章所见的制度运作落实到具体年份。其形成背景是：朝贡名义、边境安全与港口贸易的利益使交涉持续发生。 它带来的、而且仍可由后续材料追踪的改变是：它为后续册封、互市或海防安排留下文书与跨政权比较线索。 《神宗实录》万历四十六年条；《明史·外国传》；《朝鲜王朝实录》、琉球《历代宝案》或海外档案 提供相互检验的起点；但桥接条目只标示可回查的年度问题与材料链；实录有中央选择性，崇祯期尤其需要奏疏、地方、朝鲜及满文材料交叉。；本条特别核对报称信息的来源、抄传与行政分类。 因此本段仅据材料确认事件的时序、行动和可见后果，并把政策文本与地方实践、报称数字和社会经历分开处理。

\eventanchor{ML-1619-001}{明·1619（万历四十七年）四月18日}
\paragraph{1619（万历四十七年）四月18日：萨尔浒之战：明军被后金歼灭} 这项记录把本章所见的制度运作落实到具体年份。其形成背景是：前期边防、地方武装或跨区域资源调动的压力推动了该行动。 它带来的、而且仍可由后续材料追踪的改变是：它改变了后续调兵、补给或地方秩序的条件；战果和伤亡须分开核对。 《神宗实录》万历四十七年条；《明史·兵志》及相关列传；边镇、地方志或对方材料 提供相互检验的起点；但译写候选以编年材料定位；实录记载反映朝廷选择，崇祯期须另以奏疏、地方及敌对政权材料交叉。 因此本段仅据材料确认事件的时序、行动和可见后果，并把政策文本与地方实践、报称数字和社会经历分开处理。

\eventanchor{ML-1619-002}{明·1619（万历四十七年）七月26日}
\paragraph{1619（万历四十七年）七月26日：开原之战：后金攻占开原} 这项记录把本章所见的制度运作落实到具体年份。其形成背景是：前期边防、地方武装或跨区域资源调动的压力推动了该行动。 它带来的、而且仍可由后续材料追踪的改变是：它改变了后续调兵、补给或地方秩序的条件；战果和伤亡须分开核对。 《神宗实录》万历四十七年条；《明史·兵志》及相关列传；边镇、地方志或对方材料 提供相互检验的起点；但译写候选以编年材料定位；实录记载反映朝廷选择，崇祯期须另以奏疏、地方及敌对政权材料交叉。 因此本段仅据材料确认事件的时序、行动和可见后果，并把政策文本与地方实践、报称数字和社会经历分开处理。

\eventanchor{ML-1619-003}{明·1619（万历四十七年）九月3日}
\paragraph{1619（万历四十七年）九月3日：铁岭之战：后金占领铁岭} 这项记录把本章所见的制度运作落实到具体年份。其形成背景是：前期边防、地方武装或跨区域资源调动的压力推动了该行动。 它带来的、而且仍可由后续材料追踪的改变是：它改变了后续调兵、补给或地方秩序的条件；战果和伤亡须分开核对。 《神宗实录》万历四十七年条；《明史·兵志》及相关列传；边镇、地方志或对方材料 提供相互检验的起点；但译写候选以编年材料定位；实录记载反映朝廷选择，崇祯期须另以奏疏、地方及敌对政权材料交叉。 因此本段仅据材料确认事件的时序、行动和可见后果，并把政策文本与地方实践、报称数字和社会经历分开处理。

\eventanchor{ML-1619-004}{明·1619（万历四十七年）九月}
\paragraph{1619（万历四十七年）九月：首位俄罗斯特使伊万·佩特林抵达北京} 这项记录把本章所见的制度运作落实到具体年份。其形成背景是：继承、官僚协调或皇权运作的压力使问题进入朝廷程序。 它带来的、而且仍可由后续材料追踪的改变是：它影响后续人事与政策议程；动机和派系标签不能由单一叙事裁定。 《神宗实录》万历四十七年条；《明史·本纪》及相关列传；诏令、奏疏或会典版本 提供相互检验的起点；但译写候选以编年材料定位；实录记载反映朝廷选择，崇祯期须另以奏疏、地方及敌对政权材料交叉。 因此本段仅据材料确认事件的时序、行动和可见后果，并把政策文本与地方实践、报称数字和社会经历分开处理。

\subsection{1619年前后的续接}

\eventanchor{ML-1619-005}{明·1619（万历四十七年）}
\paragraph{1619（万历四十七年）：东林、书院、刻书和宗教活动的本年材料记录精英网络，社团存在不等于一致政治立场。} 这项记录把本章所见的制度运作落实到具体年份。其形成背景是：制度、教育或知识传播的需要推动了相关行动或文本形成。 它带来的、而且仍可由后续材料追踪的改变是：它提供制度和传播史的锚点，不能据此推断社会整体接受。 《神宗实录》万历四十七年条；《明史·选举志》或《艺文志》；文集、碑刻或书目材料 提供相互检验的起点；但桥接条目只标示可回查的年度问题与材料链；实录有中央选择性，崇祯期尤其需要奏疏、地方、朝鲜及满文材料交叉。 因此本段仅据材料确认事件的时序、行动和可见后果，并把政策文本与地方实践、报称数字和社会经历分开处理。

\eventanchor{ML-1620-001}{明·1620（万历四十八年）八月18日}
\paragraph{1620（万历四十八年）八月18日：万历皇帝驾崩} 这项记录把本章所见的制度运作落实到具体年份。其形成背景是：继承、官僚协调或皇权运作的压力使问题进入朝廷程序。 它带来的、而且仍可由后续材料追踪的改变是：它影响后续人事与政策议程；动机和派系标签不能由单一叙事裁定。 《神宗实录》万历四十八年条；《明史·本纪》及相关列传；诏令、奏疏或会典版本 提供相互检验的起点；但译写候选以编年材料定位；实录记载反映朝廷选择，崇祯期须另以奏疏、地方及敌对政权材料交叉。 因此本段仅据材料确认事件的时序、行动和可见后果，并把政策文本与地方实践、报称数字和社会经历分开处理。

\eventanchor{ML-1620-002}{明·1620（万历四十八年）八月28日}
\paragraph{1620（万历四十八年）八月28日：朱常洛即位太常皇帝} 这项记录把本章所见的制度运作落实到具体年份。其形成背景是：继承、官僚协调或皇权运作的压力使问题进入朝廷程序。 它带来的、而且仍可由后续材料追踪的改变是：它影响后续人事与政策议程；动机和派系标签不能由单一叙事裁定。 《神宗实录》万历四十八年条；《明史·本纪》及相关列传；诏令、奏疏或会典版本 提供相互检验的起点；但译写候选以编年材料定位；实录记载反映朝廷选择，崇祯期须另以奏疏、地方及敌对政权材料交叉。 因此本段仅据材料确认事件的时序、行动和可见后果，并把政策文本与地方实践、报称数字和社会经历分开处理。

\eventanchor{ML-1620-003}{明·1620（万历四十八年）九月6日}
\paragraph{1620（万历四十八年）九月6日：太常皇帝病了} 这项记录把本章所见的制度运作落实到具体年份。其形成背景是：继承、官僚协调或皇权运作的压力使问题进入朝廷程序。 它带来的、而且仍可由后续材料追踪的改变是：它影响后续人事与政策议程；动机和派系标签不能由单一叙事裁定。 《神宗实录》万历四十八年条；《明史·本纪》及相关列传；诏令、奏疏或会典版本 提供相互检验的起点；但译写候选以编年材料定位；实录记载反映朝廷选择，崇祯期须另以奏疏、地方及敌对政权材料交叉。 因此本段仅据材料确认事件的时序、行动和可见后果，并把政策文本与地方实践、报称数字和社会经历分开处理。

\eventanchor{ML-1620-004}{明·1620（万历四十八年）九月26日}
\paragraph{1620（万历四十八年）九月26日：太常帝驾崩} 这项记录把本章所见的制度运作落实到具体年份。其形成背景是：继承、官僚协调或皇权运作的压力使问题进入朝廷程序。 它带来的、而且仍可由后续材料追踪的改变是：它影响后续人事与政策议程；动机和派系标签不能由单一叙事裁定。 《神宗实录》万历四十八年条；《明史·本纪》及相关列传；诏令、奏疏或会典版本 提供相互检验的起点；但译写候选以编年材料定位；实录记载反映朝廷选择，崇祯期须另以奏疏、地方及敌对政权材料交叉。 因此本段仅据材料确认事件的时序、行动和可见后果，并把政策文本与地方实践、报称数字和社会经历分开处理。

\eventanchor{ML-1620-005}{明·1620（万历四十八年）十月1日}
\paragraph{1620（万历四十八年）十月1日：朱由娇成为天启皇帝} 这项记录把本章所见的制度运作落实到具体年份。其形成背景是：继承、官僚协调或皇权运作的压力使问题进入朝廷程序。 它带来的、而且仍可由后续材料追踪的改变是：它影响后续人事与政策议程；动机和派系标签不能由单一叙事裁定。 《神宗实录》万历四十八年条；《明史·本纪》及相关列传；诏令、奏疏或会典版本 提供相互检验的起点；但译写候选以编年材料定位；实录记载反映朝廷选择，崇祯期须另以奏疏、地方及敌对政权材料交叉。 因此本段仅据材料确认事件的时序、行动和可见后果，并把政策文本与地方实践、报称数字和社会经历分开处理。

\eventanchor{ML-1620-006}{明·1620（万历四十八年）}
\paragraph{1620（万历四十八年）：明代铸造厂开始生产红衣袍。} 这项记录把本章所见的制度运作落实到具体年份。其形成背景是：继承、官僚协调或皇权运作的压力使问题进入朝廷程序。 它带来的、而且仍可由后续材料追踪的改变是：它影响后续人事与政策议程；动机和派系标签不能由单一叙事裁定。 《神宗实录》万历四十八年条；《明史·本纪》及相关列传；诏令、奏疏或会典版本 提供相互检验的起点；但译写候选以编年材料定位；实录记载反映朝廷选择，崇祯期须另以奏疏、地方及敌对政权材料交叉。 因此本段仅据材料确认事件的时序、行动和可见后果，并把政策文本与地方实践、报称数字和社会经历分开处理。
